% paper/method.tex
% LaTeX outline of the GRACE Method section.
% Paste into a paper template (acmart, IEEEtran, llncs, etc.).

\section{Method}\label{sec:method}

We propose \textbf{GRACE}, a planner for embodied LLM agents that
integrates three components into a single inference pipeline:
(i) a lightweight, rule-based symbolic precondition verifier,
(ii) hierarchical decomposition with sub-goal--localised repair, and
(iii) a hybrid retrieval memory seeded with curated ground-truth plans
and grown by successful runtime episodes.

Figure~\ref{fig:overview} summarises the architecture. The remainder of
this section formalises the problem (\S\ref{sec:problem}), then describes
each component in turn (\S\ref{sec:verifier}--\S\ref{sec:memory}),
and finally presents the algorithms for initial planning and repair
(\S\ref{sec:algorithm}).

\subsection{Problem formulation}\label{sec:problem}

A household task is a tuple $\tau = (\ell, \mathcal{S}_0)$ where
$\ell$ is a natural-language description and
$\mathcal{S}_0 = (o_0, V_0)$ is the initial scene context: a textual
observation $o_0$ and a set of visible objects $V_0$. A plan
$\pi = (a_1, \dots, a_T)$ is a sequence of actions
$a_t = (\mathrm{verb}_t, \mathrm{obj}_t)$ drawn from a fixed
vocabulary $\mathcal{A}$ (Table~\ref{tab:actions} —
\texttt{MoveTo}, \texttt{Find}, \texttt{Pick}, \texttt{Place},
\texttt{Open}, \texttt{Close}, \texttt{TurnOn}, \texttt{TurnOff}, etc.).

Each action has a precondition $\mathrm{pre}(a)$ and an effect
$\mathrm{eff}(a)$ defined over a symbolic state $\sigma$ with five
fields: $\mathrm{arrived}, \mathrm{found}, \mathrm{holding},
\mathrm{opened}, \mathrm{on}$. A plan is \emph{symbolically valid}
iff $\mathrm{pre}(a_t)$ holds in
$\sigma_{t-1} = \mathrm{eff}(a_{t-1}) \circ \cdots \circ \mathrm{eff}(a_1)(\sigma_0)$
for every $t$. We let $\mathcal{V}(\pi, \sigma_0) \in \{\text{OK}\}
\cup \mathcal{R}$ denote the verifier verdict, where $\mathcal{R}$ is the
typed set of violation reasons (e.g.,
\textsc{NoFindBeforePick}, \textsc{PlaceWrongLocation}).

The planner must produce a $\pi$ that is symbolically valid and
\emph{achieves} the underlying goal when executed in the AI2-THOR
simulator. We assume access only to an LLM $f_\theta$ and the
verifier $\mathcal{V}$; we do not assume a PDDL domain or a
classical planner.

\subsection{Symbolic precondition verifier}\label{sec:verifier}

The verifier is a deterministic, hand-written function in
$\sim$250 lines of Python over the schema documented in
\texttt{pyplanner.base.STEP\_SCHEMA}. It is critical to the
contribution that the verifier \emph{never} invokes an LLM:
$\mathcal{V}$ runs in $O(|\pi|)$ time, costs zero tokens, and gives a
\emph{typed} reason for any rejection rather than the free-form
critique characteristic of self-refinement methods.

Algorithm~\ref{alg:verify} sketches \texttt{simulate}.
$\mathrm{pre}(\cdot)$ is implemented as a switch over the action's
verb; the most common rules are:
\begin{itemize}
  \item \texttt{Pick} requires $\mathrm{found} \neq \emptyset$
    and $\mathrm{holding} = \emptyset$.
  \item \texttt{Place~$r$} requires
    $\mathrm{holding} \neq \emptyset \wedge \mathrm{arrived} = r$.
  \item \texttt{Open}/\texttt{Close}/\texttt{TurnOn}/\texttt{TurnOff} require
    $\mathrm{found} = $ the action's object.
  \item Container rule: \texttt{PutIn~$c$} requires
    $c \in \mathrm{opened}$.
\end{itemize}
Rejected violations are formatted into a compact natural-language
feedback block by \texttt{format\_violations\_for\_llm}; we show in
\S\ref{sec:experiments} that this typed feedback drives more reliable
single-pass refinement than the free-form critiques used by
Self-Refine~\cite{madaan2023selfrefine}.

\begin{algorithm}
\caption{\texttt{simulate}: rule-based plan verification.}\label{alg:verify}
\begin{algorithmic}[1]
\State $\sigma \gets \sigma_0$,\ $\mathit{viol} \gets [\,]$
\For{$t = 1$ \textbf{to} $|\pi|$}
  \State $(\mathit{ok}, r) \gets \texttt{verifyStep}(a_t, \sigma)$
  \If{$\neg \mathit{ok}$}
     \State $\mathit{viol}.\textsc{append}((t, a_t, r))$;\ \textbf{break}
  \EndIf
  \State $\sigma \gets \texttt{apply}(a_t, \sigma)$
\EndFor
\State \Return $(\mathit{viol}=\emptyset,\ \mathit{viol},\ \sigma)$
\end{algorithmic}
\end{algorithm}

\subsection{Hierarchical decomposition with sub-goal blocks}\label{sec:hierarchy}

GRACE decomposes the task into 2--5 sub-goals
$g_1, \dots, g_K$ via one high-level LLM call. Each sub-goal is then
expanded by a second-level LLM call into a contiguous block of
steps $\pi^{(k)}$. We retain the sub-goal boundaries explicitly as
typed \texttt{\_SubgoalBlock} objects; the flat plan
$\pi = \pi^{(1)} \mathbin{\Vert} \cdots \mathbin{\Vert} \pi^{(K)}$ is
produced only at execution time.

Keeping the boundaries lets us perform sub-goal--localised repair
(\S\ref{sec:algorithm}). To the best of our knowledge, this is the
first work to combine hierarchical LLM planning with
\emph{suffix-only} repair scoped to the failed sub-goal; prior
hierarchical methods (e.g.~\cite{ajay2023hip},
HierarchicalFewShot in pyplanner) replan the entire remainder of the
task on failure.

\subsection{Hybrid memory retrieval}\label{sec:memory}

GRACE conditions both decomposition and expansion on $k$ retrieved
exemplars from a hybrid store $\mathcal{M} = \mathcal{M}_{\text{seed}}
\cup \mathcal{M}_{\text{live}}$:

\begin{itemize}
  \item $\mathcal{M}_{\text{seed}}$ is loaded once from a curated
    ground-truth file (38 simulator-verified plans). Seed entries
    guarantee cold-start coverage even before any episode has been
    logged. Action names are normalised to the canonical vocabulary
    via \texttt{normalize\_plan} before being stored.
  \item $\mathcal{M}_{\text{live}}$ is appended whenever the agent
    successfully completes a task. Persistence is a JSONL file;
    embedding search is optional via Chroma~\cite{chroma}, with a
    Jaccard-similarity fallback so the system runs on a minimal
    install.
\end{itemize}

Retrieval returns the top-$k$ exemplars by similarity to the task
description $\ell$. Each exemplar contributes a typed
$(\mathrm{task}, \mathrm{reasoning}, \mathrm{plan})$ triple
that the prompt assembler renders into a structured
\texttt{=== RETRIEVED EXAMPLE === } block — the same shape used by
pyplanner's \texttt{HierarchicalFewShotPlanner}, so the prompt
engineering is shared and only the \emph{contents} of the store
differ.

\subsection{Algorithms}\label{sec:algorithm}

Algorithms~\ref{alg:generate} and~\ref{alg:replan} present
$\textsc{Plan}$ and $\textsc{Replan}$ in full.
$\textsc{Plan}$ costs $1 + K + R$ LLM calls, where $K$ is the number
of sub-goals and $R \leq K$ is the number of refinement passes
triggered by verifier rejection. $\textsc{Replan}$ costs at most two
LLM calls and \emph{never} replans steps that have already been
executed.

\begin{algorithm}
\caption{\textsc{Plan} — GRACE initial planning.}\label{alg:generate}
\begin{algorithmic}[1]
\Require task $\ell$, observation $o$, visible objects $V$
\State $E \gets \textsc{Retrieve}(\mathcal{M}, \ell, k)$
\State $g_{1:K} \gets f_\theta^{\text{decompose}}(\ell, E)$
\State $\sigma \gets (\emptyset, \emptyset, \emptyset, \emptyset, \emptyset),\ V$
\State $B \gets [\,]$
\For{$k = 1$ \textbf{to} $K$}
  \State $\pi^{(k)} \gets f_\theta^{\text{expand}}(g_k, \sigma, E)$
  \For{$j = 1$ \textbf{to} $\texttt{max\_refines}$}
    \State $(\mathit{ok},\mathit{viol},\sigma') \gets \mathcal{V}(\pi^{(k)}, \sigma)$
    \If{$\mathit{ok}$}\ \textbf{break}\ \EndIf
    \State $\pi^{(k)} \gets f_\theta^{\text{refine}}(g_k, \pi^{(k)}, \mathit{viol})$
  \EndFor
  \State $\sigma \gets \texttt{rollForward}(\pi^{(k)}, \sigma)$;\
        $B.\textsc{append}((g_k, \pi^{(k)}))$
\EndFor
\State \Return $\textsc{Flatten}(B)$
\end{algorithmic}
\end{algorithm}

\begin{algorithm}
\caption{\textsc{Replan} — sub-goal--localised repair.}\label{alg:replan}
\begin{algorithmic}[1]
\Require completed steps $C$, failed step $a^\star$, reason $r$
\State $\sigma^\star \gets \texttt{simulate}(C).\mathrm{final}$
       \Comment{recover state at failure, no LLM}
\State $(g^\star, B^{>}) \gets \texttt{locate}(B, |C|)$
       \Comment{which sub-goal failed; what's still in the future}
\State $E \gets \textsc{Retrieve}(\mathcal{M}, g^\star, k)$
\State $\pi^{\text{suf}} \gets f_\theta^{\text{repair}}(g^\star,
       \sigma^\star, a^\star, r, E)$
\If{$\neg \mathcal{V}(\pi^{\text{suf}}, \sigma^\star).\mathrm{ok}$}
  \State $\pi^{\text{suf}} \gets f_\theta^{\text{refine}}(g^\star,
         \pi^{\text{suf}}, \mathit{viol})$
\EndIf
\State \Return $\pi^{\text{suf}} \mathbin{\Vert} \textsc{Flatten}(B^{>})$
\end{algorithmic}
\end{algorithm}

\paragraph{What is novel.}
Each of the three components has precedent: symbolic verification
appears in LLM+P~\cite{liu2023llmp} and ISR-LLM~\cite{zhou2024isrllm};
hierarchical decomposition in HiP~\cite{ajay2023hip}; memory-augmented
retrieval in LLM-Planner~\cite{song2023llmplanner} and
ExpeL~\cite{zhao2024expel}. GRACE's contribution is their integration
into one system under three constraints that differ from prior work:
(i) the verifier is hand-written and \emph{LLM-free}, costing zero
tokens; (ii) repair is \emph{strictly local} to the failed sub-goal,
preserving completed work and untouched future sub-goals; and
(iii) the memory store is \emph{hybrid} — both seeded from curated
ground-truth plans and grown from successful runs, with the same
prompt-block format and a Jaccard fallback when embeddings are
unavailable.
